\documentclass[12pt]{fphw}

\usepackage[utf8]{inputenc}
\usepackage[T1]{fontenc}
\usepackage{mathpazo}
\usepackage{alphabeta}
\usepackage{graphicx}
\usepackage{booktabs}
\usepackage{listings}
\usepackage{enumerate}
\usepackage{hyperref}
\usepackage{amsmath}
\usepackage{framed}
\usepackage[framed]{ntheorem}
\usepackage{amssymb}
\usepackage{xcolor}
\usepackage{tabularx}

\newframedtheorem{frm-thm}{Theorem}
\hypersetup{
    colorlinks=true,
    linkcolor=blue,
    filecolor=magenta,
    urlcolor=cyan,
}
\urlstyle{same}

\title{Deep Learning for NLP}
\author{sdi2200160}
\date{Spring Semester 2025-2026}
\institute{National and Kapodistrian University of Athens \\ Department of Informatics and Telecommunications}
\class{Artificial Intelligence II (YS19 M226 M262 M325 M405)}

\newcommand{\pending}{\textit{[fill after GPU run]}}
\newcommand{\hwonebestf}{0.45}
\newcommand{\hwonewikitwo}{0.39}
\newcommand{\hwonetwitter}{0.38}

\begin{document}
\maketitle

{
  \hypersetup{linkcolor=black}
  \tableofcontents
}
\newpage

\section{Abstract}

We study \textbf{response clarity classification} on the QEvasion political interview dataset \cite{qevasion}. Each example is a question--answer pair, and the target label is one of \emph{Clear Reply}, \emph{Ambivalent}, or \emph{Clear Non-Reply}. Homework 1 used non-contextual text representations (TF--IDF and averaged GloVe embeddings) with Logistic Regression. In Homework 2, we keep the same high-level classifier pipeline but replace the source encoder with a pretrained transformer tokenizer and replace the linear classifier with a transformer sequence classifier fine-tuned through an explicit PyTorch training loop.

Our goal in this homework is not aggressive leaderboard optimization. Instead, we adopt a \textbf{simple and principled} training recipe that is easy to justify: paired question--answer encoding, one fixed stratified validation split, AdamW optimization, class-weighted loss for imbalance, gradient clipping, and model selection by validation macro F1. We compare the three required pretrained checkpoints---BERT-base \cite{bert}, DistilBERT \cite{distilbert}, and DeBERTa-v3-base \cite{deberta}. The exact run-specific scores are inserted after executing the GPU notebooks, while the methodological comparison and the HW1 baseline reference are documented here.

\section{Data Processing and Analysis}

\subsection{Dataset}

We use the QEvasion dataset from Hugging Face \cite{qevasion}. Each supervised instance contains:
\begin{itemize}
	\item a single segmented journalist \texttt{question},
	\item the corresponding \texttt{interview\_answer},
	\item the three-class target \texttt{clarity\_label}.
\end{itemize}

As in Homework 1, we use the official train/test split provided by the dataset. The training split is used for fitting and internal validation; the official test split is reserved for final reporting and Kaggle submission generation.

\subsection{Preprocessing}

Preprocessing in Homework 2 is intentionally lighter than in Homework 1. We do \emph{not} apply manual lemmatization, stopword removal, or handcrafted punctuation normalization. Instead, we keep the raw question and answer text, strip only missing values and outer whitespace, and rely on the pretrained tokenizer to perform subword segmentation and insert the appropriate special tokens.

This design is principled for transformer fine-tuning: the checkpoints were pretrained on raw text, so aggressive manual normalization is less clearly beneficial than it was for sparse bag-of-words models. The notebook explicitly contrasts this decision with the heavier preprocessing used in Homework 1.

\subsection{Exploratory Analysis}

The exploratory analysis remains intentionally compact. We examine:
\begin{itemize}
	\item class imbalance in the train and test splits,
	\item question length, answer length, and combined length distributions,
	\item subgroup behavior for short, medium, and long question/answer buckets.
\end{itemize}

As already observed in Homework 1, the dataset is imbalanced: \emph{Clear Reply} is the majority class, while \emph{Ambivalent} is the smallest and most difficult class. This motivates class-weighted training and macro-averaged evaluation.

\subsection{Input Construction}

Unlike Homework 1, where we concatenated the question and answer into one string before feature extraction, Homework 2 preserves them as a \textbf{sentence pair}. The tokenizer receives two sequences:
\[
(\text{question}, \text{answer})
\]
and converts them into the checkpoint-specific sequence classification format. Conceptually, the model input is:
\[
[\text{CLS}]~\text{question}~[\text{SEP}]~\text{answer}~[\text{SEP}]
\]
for BERT-like architectures, although the exact special tokens are handled internally by each tokenizer.

\section{Models and Training Setup}

\subsection{Shared Pipeline}

The abstract pipeline is unchanged from Homework 1:
\begin{enumerate}
	\item preprocess the decoded input,
	\item encode the source text,
	\item encode the target labels,
	\item fit the model,
	\item decode predictions and evaluate with macro-averaged metrics.
\end{enumerate}

We keep the same lightweight \texttt{Classifier} wrapper. This is pedagogically useful because Homework 2 becomes a direct extension of Homework 1 rather than a completely different codebase.

\subsection{Transformer Models}

We compare the three required checkpoints:
\begin{itemize}
	\item \texttt{bert-base-uncased} \cite{bert},
	\item \texttt{distilbert-base-uncased} \cite{distilbert},
	\item \texttt{microsoft/deberta-v3-base} \cite{deberta}.
\end{itemize}

All models are fine-tuned as three-way sequence classifiers using PyTorch and Hugging Face Transformers \cite{huggingface}. We do \emph{not} use the Hugging Face Trainer API, in accordance with the assignment specification; the training loop is implemented explicitly.

\subsection{Validation Strategy}

Homework 1 used grid search with 3-fold cross-validation because the models were classical and cheap to train. Homework 2 uses transformer fine-tuning, so repeating full training three times per hyperparameter setting would be unnecessarily expensive. We therefore adopt a single \textbf{15\% stratified validation split} drawn from the official training split:
\begin{itemize}
	\item 85\% of the training split is used for gradient updates,
	\item 15\% is used for validation during training,
	\item the official test split is never used for model selection.
\end{itemize}

This choice keeps the procedure computationally reasonable while still giving us a stable validation signal for checkpoint selection.

\subsection{Training Recipe}

Our shared default recipe is intentionally small and consistent across models:
\begin{itemize}
	\item optimizer: AdamW \cite{adamw},
	\item learning rate: $2 \times 10^{-5}$,
	\item batch size: 16,
	\item weight decay: 0.01,
	\item epochs: 4,
	\item maximum sequence length: 256,
	\item gradient clipping: 1.0,
	\item class weighting: enabled,
	\item random seed: 42.
\end{itemize}

The only sanctioned tuning is a \textbf{small uniform sanity sweep} applied separately to each required model:
\begin{itemize}
	\item epochs $\in \{3, 4\}$,
	\item maximum length $\in \{128, 256, 384\}$.
\end{itemize}

This produces 6 runs per model, or 18 runs in total across BERT, DistilBERT, and DeBERTa. The same search space is used for all three checkpoints. This keeps the tuning budget uniform while still remaining small enough to fit the pedagogical scope of the assignment.

\begin{table}[h]
\centering
\small
\begin{tabularx}{\textwidth}{Xlll}
\toprule
Model & Validation Strategy & Tuning & Final Recipe \\
\midrule
BERT & 15\% stratified split & same 2 x 3 sweep & \pending \\
DistilBERT & 15\% stratified split & same 2 x 3 sweep & \pending \\
DeBERTa-v3 & 15\% stratified split & same 2 x 3 sweep & \pending \\
\bottomrule
\end{tabularx}
\caption{Training protocol for the three required transformer models.}
\label{tab:training-protocol}
\end{table}

\section{Evaluation Protocol}

We report the same metrics as in Homework 1:
\begin{itemize}
	\item Accuracy,
	\item Macro Precision,
	\item Macro Recall,
	\item Macro F1.
\end{itemize}

Macro F1 remains the primary metric because it weights all classes equally and therefore reflects performance on the minority \emph{Ambivalent} class more fairly than accuracy alone.

In addition to overall scores, the notebook also produces:
\begin{itemize}
	\item confusion matrices,
	\item learning curves (train loss, validation loss, validation macro F1),
	\item subgroup tables by question length and answer length,
	\item concrete misclassified examples for error analysis.
\end{itemize}

\section{Results and Comparison}

\subsection{Reference Point from Homework 1}

We do not rerun the classical baselines inside Homework 2. Instead, we reference the best values already obtained in Homework 1:

\begin{table}[h]
\centering
\begin{tabular}{ll}
\toprule
Homework 1 Model & Macro F1 (reference) \\
\midrule
TF--IDF + Logistic Regression & \hwonebestf \\
GloVe Wiki + Logistic Regression & \hwonewikitwo \\
GloVe Twitter + Logistic Regression & \hwonetwitter \\
\bottomrule
\end{tabular}
\caption{Homework 1 reference results, taken from the previous report.}
\label{tab:hw1-reference}
\end{table}

These values are important because they provide the exact comparison point that motivates Homework 2: if contextual transformers are working as intended, they should improve on averaged word embeddings and ideally also exceed the strong TF--IDF baseline.

\subsection{Homework 2 Model Comparison}

\begin{table}[h]
\centering
\small
\begin{tabularx}{\textwidth}{Xcccc}
\toprule
Model & Accuracy & Macro Precision & Macro Recall & Macro F1 \\
\midrule
BERT & \pending & \pending & \pending & \pending \\
DistilBERT & \pending & \pending & \pending & \pending \\
DeBERTa-v3 & \pending & \pending & \pending & \pending \\
\bottomrule
\end{tabularx}
\caption{Homework 2 transformer results on the official test split.}
\label{tab:hw2-results}
\end{table}

\subsection{Interpretation}

The main hypothesis is that transformer models improve on Homework 1 because they encode:
\begin{itemize}
	\item word order,
	\item contextual meaning,
	\item the interaction between the journalist's question and the interviewee's answer.
\end{itemize}

These are exactly the aspects that TF--IDF and mean-pooled GloVe struggle to represent well. In particular, the clarity task depends not only on what words appear in the answer, but also on whether the answer is responsive to the specific question asked. Sentence-pair transformers are a natural fit for this interaction.

After notebook execution, Table \ref{tab:hw2-results} should be read together with Table \ref{tab:hw1-reference}. The comparison should answer two questions:
\begin{enumerate}
	\item Does contextual fine-tuning clearly outperform the Homework 1 GloVe baselines?
	\item Does it also beat the stronger TF--IDF baseline?
\end{enumerate}

\section{Subgroup Analysis}

The assignment explicitly requests analysis across data subgroups. We therefore evaluate the final predictions on buckets formed by question length and answer length:
\begin{itemize}
	\item short,
	\item medium,
	\item long.
\end{itemize}

\begin{table}[h]
\centering
\small
\begin{tabularx}{\textwidth}{Xccc}
\toprule
Bucket & Count & Accuracy & Macro F1 \\
\midrule
Short & \pending & \pending & \pending \\
Medium & \pending & \pending & \pending \\
Long & \pending & \pending & \pending \\
\bottomrule
\end{tabularx}
\caption{Performance across question-length buckets.}
\label{tab:question-subgroups}
\end{table}

\begin{table}[h]
\centering
\small
\begin{tabularx}{\textwidth}{Xccc}
\toprule
Bucket & Count & Accuracy & Macro F1 \\
\midrule
Short & \pending & \pending & \pending \\
Medium & \pending & \pending & \pending \\
Long & \pending & \pending & \pending \\
\bottomrule
\end{tabularx}
\caption{Performance across answer-length buckets.}
\label{tab:answer-subgroups}
\end{table}

These subgroup tables are useful for testing whether the models struggle more with long, potentially more evasive answers, or with short answers that lack lexical evidence.

\section{Error Analysis}

Error analysis is a central part of this assignment. The notebook extracts the most common confusion pairs and several representative mistakes. The report should summarize the patterns that emerge after execution. The following hypotheses are especially relevant:
\begin{enumerate}
	\item \textbf{Ambivalent} is likely to remain the hardest class.
	\item Clear Non-Reply and Ambivalent may still be confused when the answer is topically related but not directly responsive.
	\item Longer answers may hurt performance when they contain partial responsiveness mixed with rhetorical evasion.
\end{enumerate}

Concrete examples should be inserted here after running the notebooks. A useful structure is:
\begin{itemize}
	\item one example where the model predicts Clear Reply instead of Ambivalent,
	\item one example where it predicts Ambivalent instead of Clear Non-Reply,
\item one example where contextual modeling helps relative to the Homework 1 baselines.
\end{itemize}

\section{Discussion}

This simplified Homework 2 design deliberately avoids heavy optimization. That is a feature, not a limitation: it makes the core conceptual improvement over Homework 1 easier to understand. The important difference is not a large hyperparameter search, but the move from non-contextual features to a contextual sentence-pair encoder.

The design choices can be summarized as follows:
\begin{itemize}
	\item use raw question--answer pairs rather than handcrafted text normalization,
	\item fine-tune standard pretrained checkpoints instead of building custom architectures,
	\item use one fixed validation split rather than expensive cross-validation,
	\item keep the same pipeline abstraction as Homework 1 for clarity of comparison.
\end{itemize}

This setup makes it easy to argue why the transformer approach should improve: the model sees the full question--answer interaction and can condition the meaning of each token on the surrounding context.

\section{Conclusion}

Homework 2 extends Homework 1 in the most direct possible way: we preserve the pipeline abstraction and replace only the representation and classifier components. This makes the comparison between the two homeworks clean and interpretable. Homework 1 established that TF--IDF was much stronger than averaged GloVe on this task. Homework 2 tests whether contextual fine-tuning can move beyond both of those non-contextual baselines.

After the GPU notebooks are executed, the remaining task is to populate the result tables, insert the subgroup values, and summarize the concrete failure cases found in the error-analysis cells. The underlying methodology is already fixed: a minimal transformer training loop, a single stratified validation split, the same small 2 x 3 sweep for each required model, and a direct comparison to the Homework 1 reference results.

\nocite{*}
\bibliographystyle{plain}
\bibliography{refs}

\end{document}
